\section{Conclusions}\label{sec:concl}
This paper has studied the integrated train timetabling, stop planning, and platforming problem with the timetable makespan as the primary objective. We represented the problem on an event-activity network whose activities between trains are disjunctive and depend on the stopping and station routing decisions, and formulated it as a MILP with continuous event times. For a fixed structure of stops, orders, and station routings, the minimum makespan equals the length of a longest path in the selected network, and the station platforming problems decompose by station. Building on these properties, we developed a conflict-driven ALNS with nine destroy operators. Its joint neighborhoods couple the resolution of station conflicts with the coordinated changes to stops, arrival and departure times, and station routings that are needed to keep every affected service count within its bounds, and a restricted MILP repairs them while retaining all constraints that link free and fixed decisions.

In the computational experiments, CD-ALNS reaches the proven optimal makespan in the small-scale case. In the medium- and large-scale cases, it quickly finds jointly feasible timetables and further reduces their makespans, whereas direct MILP solution and an adapted LBBD find no jointly feasible timetable within six hours. In the paired medium-scale comparison, the macroscopic ALNS phase yields earlier joint feasibility and earlier attainment of the same final makespan, with worse secondary objectives. \tbd~\textit{Summarize here the results of the experiments E1--E5.}

Several directions remain for future research. First, lower bounds for the medium- and large-scale cases would allow the optimality gaps to be evaluated. Second, the approach could be extended to bidirectional lines with single-track sections and to network timetables with connections. Third, passenger-oriented objectives and robustness against delays could be added to the lexicographic objective, and the neighborhood construction could be combined with learning-based operator selection.
